\section{Introduction}\label{sec:introduction}

Strong cosmic censorship asks whether generic initial data determine a
spacetime that cannot be continued beyond its maximal globally hyperbolic
development. The regularity of the proposed continuation is part of the
question. Here we consider \(C^2\) future extensions of vacuum developments
near the complete two-ended bridge of an arbitrary fixed rotating
subextremal Kerr black hole. No symmetry is imposed on the perturbation.

The exact Kerr development has a smooth Cauchy horizon. Nevertheless,
we prove that data admitting any \(C^2\) future extension are meagre in
a weighted smooth neighborhood of its bridge data. The neighborhood
requires smallness at only one finite order, and the allowed tails need
not have prescribed leading coefficients. The conclusion concerns the
full development, not just the portion covered by an exterior or
double-null coordinate construction.

\subsection{The precise statement}

Fix \(M>0\) and \(0<\mathfrak a<M\). Let
\((\Sigma,h_*,K_*)\), \(\Sigma\simeq\mathbb R\times\mathbb S^2\),
be the complete two-ended Kerr datum through the bifurcation sphere,
with its chosen future orientation. Fix the connection \(\nabla_*\)
and tensor norm \(|\cdot|_*\) of \(h_*\), and a smooth
\(\rho\geq1\) agreeing with the Kerr radial coordinate outside a compact
set on each end. For a pair of smooth symmetric covariant two-tensors,
define
\begin{equation}\label{intro:seminorms}
 p_m(u,v)=\sum_{j=0}^{m}\left[
 \sup_\Sigma\rho^{1+j}|\nabla_*^j u|_*
 +\sup_\Sigma\rho^{2+j}|\nabla_*^j v|_*\right],
 \qquad m=0,1,2,\ldots.
\end{equation}
Let \(E\) be the space on which every \(p_m\) is finite. Let
\(\mathcal D\subset(h_*,K_*)+E\) consist of the smooth vacuum data
\(d=(h,K)\) with \(h\) positive definite and complete, satisfying
\begin{equation}\label{intro:constraints}
 R(h)+(\operatorname{tr}_hK)^2-|K|_h^2=0,
 \qquad \operatorname{div}_hK-\mathrm d(\operatorname{tr}_hK)=0.
\end{equation}
We give \(\mathcal D\) the relative topology of all the seminorms
\eqref{intro:seminorms}. In particular, no leading coefficient in either
asymptotic tail is fixed.

Write \((M_d,g_d)\) for the full maximal globally hyperbolic vacuum
development of \(d\), with the specified Cauchy embedding of \(\Sigma\).
We use smooth vacuum Cauchy theory and geometric uniqueness in the sense
of \cite{ChoquetBruhat1952,ChoquetBruhatGeroch1969}.

\begin{definition}[Future extension]\label{intro:extension}
A \(C^2\) future extension of \((M_d,g_d)\) is a connected time-oriented
smooth four-manifold \(\widetilde M\) with a nondegenerate \(C^2\)
Lorentzian metric \(\widetilde g\), and a smooth time-orientation-preserving
isometric embedding \(\iota:M_d\longrightarrow\widetilde M\) onto a
proper open subset, such that a future-directed timelike \(C^1\) curve
\(c:[0,1]\longrightarrow\widetilde M\) satisfies
\[
 c([0,1))\subset\iota(M_d),\qquad
 c(1)\in\partial\iota(M_d).
\]
\end{definition}
We use signature \((-+++ )\). The extension need not be vacuum or
globally hyperbolic outside the image. In particular, a regular geodesic
exit is not assumed in this definition; it will follow locally from
the extension and global hyperbolicity of the original development.

\begin{theorem}[Local \(C^2\) strong cosmic censorship]\label{thm:main}
For every fixed \(M>0\) and \(0<\mathfrak a<M\), there exists
\(\epsilon>0\) such that the data in
\[
 \mathcal U_\epsilon
 =\{d\in\mathcal D:p_{10}(d-(h_*,K_*))<\epsilon\}
\]
whose full maximal globally hyperbolic vacuum developments admit a
\(C^2\) future extension form a meagre subset of \(\mathcal U_\epsilon\).
Equivalently, the future-inextendible data contain a countable intersection
of relatively open dense subsets of \(\mathcal U_\epsilon\).
\end{theorem}

Thus \Cref{thm:main} resolves positively the stated local \(C^2\)
strong cosmic censorship problem at every strictly subextremal rotating
Kerr bridge. It is not a global assertion for all vacuum initial data,
nor a claim at a lower extension regularity. Constants and
\(\epsilon\) may depend on the fixed center; no estimate uniform as
\(\mathfrak a\to0\) or \(\mathfrak a\to M\) is claimed. Although
all higher seminorms are finite, none beyond \(p_{10}\) must be small.

\subsection{History and relation to earlier work}
The Cauchy theory of Four\`es-Bruhat and Choquet-Bruhat--Geroch
\cite{ChoquetBruhat1952,ChoquetBruhatGeroch1969} makes the maximal
globally hyperbolic development the natural spacetime determined by
vacuum initial data. Sbierski \cite{Sbierski2016} gives a modern
construction with explicit initial-surface and embedding hypotheses.
The Kerr solutions \cite{Kerr1963} show why maximality in this class
need not mean inextendibility as a Lorentzian manifold: their analytic
continuation has inner Cauchy horizons, whose global geometry was
analyzed by Carter \cite{Carter1968}. Penrose's cosmic-censorship
programme began with the visibility of singularities in collapse
\cite{Penrose1969}; his later strong formulation linked Cauchy-horizon
instability to global predictability \cite[Section~12.3.2]{Penrose1979}.
The regularity of an extension and the topology used to define
genericity are essential parts of a precise formulation.

The physical blueshift and mass-inflation analyses of Poisson--Israel
and Ori \cite{PoissonIsrael1990,Ori1991,Ori1992} suggested that a
Cauchy horizon can retain a continuous metric while developing a
singularity in its derivatives. Dafermos
\cite{Dafermos2003,Dafermos2005} established rigorous nonlinear
stability and instability results for the spherically symmetric
Einstein--Maxwell--scalar-field equations. In that model,
Luk--Oh \cite{LukOh2019I,LukOh2019II} proved generic \(C^2\)
future inextendibility for admissible complete two-ended asymptotically
flat data, combining exterior lower bounds with interior instability.
Their data space and genericity topologies differ from those used
here. For the vacuum equations without symmetry, Luk
\cite{Luk2018} constructed weak null singularities, and
Dafermos--Luk \cite{DafermosLuk2017} proved continuous extendibility
across a nontrivial Cauchy-horizon portion under quantitative assumptions
on data already inside a black hole. This distinction between metric
continuity and higher regularity is central to the present problem.

Recent results strengthen the vacuum interior instability picture.
Gurriaran \cite{Gurriaran2026} proves curvature blow-up and locally
Lipschitz inextendibility across a horizon portion near timelike
infinity, assuming nonlinear Price-law asymptotics with a nonzero
nonaxisymmetric leading curvature amplitude. Luk--Sbierski
\cite{LukSbierski2026} prove formation of a continuously extendible
but non-Lipschitz weak null singularity from characteristic data with
quantitative decay and a weighted lower bound on a radiative curvature
mode. Sbierski's geometric criteria
\cite{Sbierski2022,Sbierski2026} explain how controlled parallel
transport and robust signed curvature information obstruct Lipschitz
extension. The Kerr interior instability results just cited concern a
prescribed interior or horizon portion under explicit radiation hypotheses. The problem addressed
here is to start with the complete weighted smooth bridge data and
rule out every \(C^2\) future exit generically. The finite perturbations
below provide the needed curvature test violations without imposing a
radiation lower bound on each starting datum.

\paragraph{Exterior evolution.}
The exterior estimates draw on both energy and stationary methods.
The redshift estimate and the physical-space decay method of
Dafermos--Rodnianski
\cite{DafermosRodnianski2009,DafermosRodnianski2010} were followed by
boundedness, decay, and scattering for scalar waves throughout the
subextremal Kerr range
\cite{DafermosRodnianskiShlapentokhRothman2016,DafermosRodnianskiShlapentokhRothman2018}.
For radiative spin-two fields, the Teukolsky estimates of
Dafermos--Holzegel--Rodnianski \cite{DafermosHolzegelRodnianski2019}
and Shlapentokh-Rothman--Teixeira da Costa
\cite{ShlapentokhRothmanTeixeiraDaCosta2023} treat slowly rotating and
full-subextremal Kerr, respectively. At the metric level,
H\"afner--Hintz--Vasy \cite{HafnerHintzVasy2021,HHV2025} establish
linear stability, using stationary spectral analysis and constraint
damping; the full-subextremal mode analysis by
Andersson--H\"afner--Whiting \cite{AnderssonHafnerWhiting2026}
is an important ingredient. These linear results distinguish radiative
decay from the stationary Kerr and gauge modes that must be tracked
in a nonlinear argument.

Nonlinear stability for slowly rotating Kerr was established through
the geometric and wave-equation analysis of Klainerman, Szeftel,
Giorgi, and Shen; see
\cite{KlainermanSzeftel2021,GiorgiKlainermanSzeftel2022} and the
constituent works cited there. Hintz \cite{Hintz2026} establishes
nonlinear exterior stability throughout the subextremal range for
compact perturbations and a broader partially polyhomogeneous
asymptotic class. That class has metric perturbations decaying like
\(r^{-1-\epsilon_0}\), with \(\epsilon_0>0\), and a faster-decaying
remainder; it does not directly include the arbitrary order-\(r^{-1}\)
symbol tails permitted here. We therefore prove the compact and far
estimates for our data class together. Their wave-coordinate
cancellation is related to the weak-null structure identified and
used by Lindblad--Rodnianski
\cite{LindbladRodnianski2003,LindbladRodnianski2010}.
The use of elliptic control at selected frequencies together with
frequency-localized physical-space energy currents also has a close methodological
parallel in the full-subextremal quasilinear-wave analysis of
Dafermos--Holzegel--Rodnianski--Taylor
\cite{DafermosHolzegelRodnianskiTaylor2024}. The tensorial gauge
construction and its estimates are established here for the
Einstein system.

\paragraph{Exact data and finite signals.}
Localized deformation of the constraints was developed by Corvino,
Corvino--Schoen, and Chru\'sciel--Delay
\cite{Corvino2000,CorvinoSchoen2006,ChruscielDelay2003};
Carlotto--Schoen \cite{CarlottoSchoen2016} proved a further conical
localization of asymptotically flat data. Our correction fixes an inner
region while allowing an asymptotic tail, so it can change the charges.
The explicit inverse uses a localized divergence primitive in the
tradition of Bogovskii \cite{Bogovskii1980}, whose support and regularity
properties are treated by Costabel--McIntosh
\cite{CostabelMcIntosh2010}. We prove the required inverse here because its
support, derivative gain, and nonlinear dependence are needed in the
finite experiment. Similarly, complex Gaussian beams and their
superpositions provide the linear antecedents of the oscillatory
construction \cite{Ralston1982,Tanushev2008,LiuRunborgTanushev2013}.
The matching of Einstein wave-gauge data and the bounds on time
intervals growing with the experiment parameter are proved below.


\subsection{From regular exits to finite observations}

The proof turns a possible extension into a finite condition that can
be violated by a small vacuum perturbation. A \(C^2\) extension supplies
an open family of timelike geodesics whose proper durations and curvature
components in parallel frames are bounded. A countable basis of initial
geodesic states and integer bounds give countably many such tests.
Each test is closed: a strict violation occurs on a compact geodesic
segment and persists under a sufficiently small change of the data.
Consequently, to prove meagreness, it suffices to violate the relevant
test arbitrarily close to every datum admitting a regular exit.
\Cref{sec:category} makes this reduction precise. Its decisive feature
is that one large tidal measurement at a finite point suffices; the
perturbation need not establish an asymptotic radiation law for the
starting datum.

We obtain that measurement from the blueshift of a small oscillatory
signal approaching an interior branch. In the optical coordinates used
below, a branch means a limiting side where one coordinate remains finite
and the other tends to infinity; at a corner both tend to infinity.
These are coordinate-limit regimes of the comparison region.
The main analytic work has two
parts. First we construct a comparison development and control its
geometry throughout one fixed neighborhood of Kerr data. Then, for a
particular datum and a prescribed finite approximation order, we place
a late signal in that geometry and realize it in the full development
of nearby exact vacuum data. The distinction matters: the time and
frequency of the experiment may depend on the chosen finite order,
whereas the original \(p_{10}\) neighborhood must serve all tests.

\paragraph{Exterior control and the ergoregion.}
The exterior construction combines stationary estimates in a bounded
region with energy estimates on the asymptotically flat ends.
Stationary time translation need not be timelike near all trapped rays
at arbitrary spin. We instead separate directions in the cotangent
bundle. \Cref{sec:geometry} constructs a horizon-crossing spacelike
slicing for which the spatial projections of the trapped null
covectors lie in a cone where the zero-temporal-frequency stationary
symbol is elliptic. This directional ellipticity persists even when
the base point lies in the ergoregion. \Cref{sec:stationary} uses the
cone to localize the additional gauge terms and calibrate tensor
transport while preserving the stationary modes and their
time-translation chains. The stationary analysis of
\cite{HHV2025} and the normally hyperbolic propagation estimates of
\cite{Dyatlov2016,Hintz2017} enter at the points specified there.

The compact estimate in \Cref{sec:compact} couples the two-derivative
elliptic gain on that cone to a one-derivative propagation estimate,
with one additional source derivative, away from the projected trapped
sets. On the ends, the weak-null cancellations and weighted energy
currents control the arbitrary symbol-bounded tails in
\eqref{intro:seminorms}. \Cref{sec:far} closes these far estimates
simultaneously with the compact estimates and the evolution of the
Kerr parameters. The detailed far energy and source arguments are in
\Cref{fe:section,fs:section}. All of these are estimates on an arbitrary
finite interval of smooth existence, with bounds that permit
continuation; a global nonlinear stability statement for this data
class is not assumed. Once the low-order bounds close, each fixed
higher derivative order has a subexponential bound with constants
depending on the smooth datum.

\paragraph{The interior and every possible exit.}
\Cref{sec:interior} prepares the data entering the black hole and
constructs a double-null comparison region. Exact Kerr optical
profiles, with their two ends normalized to a common positive inner
rate, organize the lapse, angular metric, and transport equations.
The low-order estimates control the geometry; the higher estimates
control the finite derivatives needed later by the oscillatory
construction. \Cref{sec:comparison} joins this region to the exterior
and proves that the entire original bridge is Cauchy. Geometric
uniqueness therefore embeds the resulting comparison development in
the full maximal development of the same data.

We then follow the geodesic family supplied by a putative extension.
Exterior completeness and the causal geometry force its relevant
first exit from the comparison region to occur at an interior branch
or in the corner regime. In the corner case, a small
change of initial geodesic state within the same open family produces
runout at a noncorner branch of the comparison region. Thus the original
closed curvature test still applies
to the selected geodesic. This step connects the local interior
coordinates to an arbitrary future extension of the full development;
it does not require the comparison region to exhaust that development.

\paragraph{A small exact vacuum perturbation with a large tidal effect.}
\Cref{sec:packet} constructs a high-frequency signal near a late point
of the selected branch. The branch geodesic measures an exponentially
increasing optical frequency. Its contraction with the oscillatory
curvature therefore yields a large tidal component even when the
signal is very small at entry. A return-ray calculation and a complex
Gaussian beam carry the required finite entry derivatives backwards
to a distant initial annulus, with controlled losses. The subsequent
forward comparison uses ordinary local wave-map reduction with the
actual background metric as target.

An approximate signal does not yet give vacuum initial data.
The exterior solve proved in \Cref{pk:constraints} and applied in
\Cref{pk2:constraints} supplies an exact correction that vanishes on a
protected inner ball and retains a decaying tail at infinity. The tail
allows the asymptotic charges to change and is included in the complete
corrected bridge. A causal estimate shows that this correction cannot
alter the earlier entry segment on which the comparison relies.
Finally, the finite Cauchy attachment in \Cref{sec:packet} places the
tidal observation, its frame, and its geodesic back to the original
initial surface in the full development of the corrected data.
The resulting strict violation makes the corresponding extendible
set nowhere dense. The countable union from \Cref{sec:category}
completes the proof.

Figure~\ref{fig:proof} displays the two constructions and the point
at which they meet. The fixed background neighborhood is chosen before
the test or approximation order. The finite derivative orders and
their losses are chosen before the small frequency exponent; the
formal accuracy and late time are then chosen to achieve the specified
tolerance. In particular, the construction requires neither smallness
of all higher seminorms nor one perturbation sequence converging in
every seminorm at once.

\begin{figure}[htbp]
\centering
\begin{tikzpicture}[
 box/.style={draw,align=center,text width=4.8cm,inner sep=7pt,font=\small},
 arr/.style={-{Stealth[length=2mm]},thick},node distance=1.05cm and 1.3cm]
 \node[box] (base) {One small \(p_{10}\) neighborhood\\all higher seminorms finite};
 \node[box,below=of base] (comparison) {Background comparison\\development\\complete original bridge\\Cauchy};
 \node[box,below=of comparison] (branch) {Branch seed in the original open family\\closed test bounds parallel curvature};
 \node[box,right=of branch] (signal) {Late finite tidal experiment\\parameters depend on a finite order};
 \node[box,above=of signal] (corrected) {Exact corrected vacuum data\\arbitrarily close in that order};
 \node[box,above=of corrected] (mg) {Finite Cauchy attachment\\observation in the new full development};
 \draw[arr] (base)--(comparison);
 \draw[arr] (comparison)--(branch);
 \draw[arr] (branch)--(signal);
 \draw[arr] (signal)--(corrected);
 \draw[arr] (corrected)--(mg);
\end{tikzpicture}
\caption{The comparison region locates a possible regular exit. The
separate finite experiment violates its closed curvature test in the
full development of nearby data. Neither an exhaustive description
of that development nor smallness at every derivative order is needed.}
\label{fig:proof}
\end{figure}

The intermediate results also separate the evolution problem from
the regularity of a proposed extension. The complete comparison
development, the entry and normalized double-null estimates, the
linear beam preparation, the exterior constraint inverse, and the
whole-data attachment can be used when the final obstruction is
holonomy rather than curvature. For \(C^1\) extensions and for
continuous extensions with locally square-integrable connection,
one must additionally control signed or averaged parallel transport.
Those uses require their own nonlinear amplitude estimates: the
tidal normalization in this paper does not supply the larger-amplitude
comparisons needed for these weaker regularities.

These regularity refinements are established in
\cite[Theorem~2.3]{KerrCOne} and
\cite[Theorem~1.2]{KerrWeakConnection}, respectively, as local residual
statements near each fixed rotating subextremal Kerr bridge.

\subsection{Scaling and conventions}

It suffices to prove the theorem with \(M=1\). Under the constant
rescaling \(g\mapsto M^{-2}g\), the induced covariant data become
\((M^{-2}h,M^{-1}K)\). After identifying the bridge by dimensionless
Kerr coordinates, this is a homeomorphism of the constraint spaces;
the background connections agree and the tensor norms and radial
weights change by fixed-order comparable factors. Completeness,
vacuum maximality and the extension property are preserved. An
original \(p_{10}\) ball can therefore be chosen inside the inverse
image of the rescaled neighborhood, and relative meagreness restricts
to that open ball. We retain \(M\) in geometric identities when doing
so clarifies their dependence on the background.

We write \(\mathfrak a\) for Kerr rotation when another local quantity
is called \(a\), notably a parameter derivative in exterior estimates
or a null lapse coefficient in the interior. All spacetime and
angular Sobolev norms are computed in the specified fixed charts or
frames; on fixed supports, smooth changes give equivalent norms.
Supported and extendible spaces at an outflow boundary are dual in
the usual way. Symbol bounds on an end include all differentiated
bounds, with an additional inverse radial power for each unscaled
spatial derivative. A bound \(e^{o(T)}\) means that for every
\(\sigma>0\) the quantity is at most \(C_\sigma e^{\sigma T}\);
the constant may depend on the smooth datum and on a fixed derivative
order. These conventions will always be used at fixed subextremal
parameters.
